\chapter{Минимальный ковариантный Стайнспринг-носитель}
\label{ch:v8-minimal-covariant-stinespring-carrier-gate}

\section{Из генератора в точный одношаговый канал}

Пусть $D_a$, $a=1,\ldots,12$, --- самосопряжённые операторы полного
межсекторных стрелок базиса на $\mathbb C^{11}\oplus\mathbb C^{10}$ и
\begin{equation}
 G=\sum_{a=1}^{12}D_a^2.
\end{equation}
Для $p\geq0$ определим Краус-операторы
\begin{equation}
 K_0(p)=\sqrt{I-pG},
 \qquad
 K_a(p)=\sqrt p\,D_a.
 \label{eq:v8-one-step-stinespring-kraus}
\end{equation}
Поскольку
\begin{equation}
 \Spec G=\{0^{(9)},1^{(6)},2^{(3)},3^{(2)},6^{(1)}\},
\end{equation}
формула положительна точно в интервале
\begin{equation}
 0\leq p\leq\frac16.
 \label{eq:v8-stinespring-step-window}
\end{equation}
Полученная карта
\begin{equation}
 \Phi_p(X)=\sum_{a=0}^{12}K_a(p)^*XK_a(p)
 \label{eq:v8-exact-one-step-channel}
\end{equation}
точно унитальна, сохраняет след и сохраняет концевую алгебру
$M_{11}(\mathbb C)\oplus M_{10}(\mathbb C)$.

\section{Минимальность и происхождение среды}

При любом внутреннем $0<p<1/6$ тринадцать операторов
$K_0,K_1,\ldots,K_{12}$ линейно независимы. Поэтому Чой/Краус-ранг карты
равен $13$, и минимальный Стайнспринг-носитель имеет вид
\begin{equation}
 \mathcal K_{\rm env}
 =\mathbb C|0\rangle\oplus
 \bigl(E_{q\ell}^{\rm cross}\otimes_{\mathbb R}\mathbb C\bigr),
 \qquad \dim_{\mathbb C}\mathcal K_{\rm env}=13,
 \quad \dim_{\mathbb R}E_{q\ell}^{\rm cross}=12.
 \label{eq:v8-minimal-environment-carrier}
\end{equation}
Нетривиальная часть среды не является новым физическим мультиплетом: это
комплексификация пространства меток уже существующей вещественной
межсекторных стрелок опоры, а не удвоенный набор физических полей. Дополнительная линия
$|0\rangle$ несёт тривиальное представление и кодирует отсутствие скачка.

При случайных преобразованиях $SU(3)\times SU(2)\times U(1)$ двенадцать
$D_a$ смешиваются ортогональной вещественной матрицей. Максимальные ошибки
ортогональности, восстановления операторов и ковариантности канала меньше
$6\cdot10^{-15}$. Следовательно, среда несёт не произвольно назначенное, а
уже известное представление пространства стрелок.

\section{Граница одношагового результата}

Производная при $p=0$ возвращает требуемый генератор:
\begin{equation}
 \left.\frac{d\Phi_p}{dp}\right|_{p=0}(X)
 =\sum_aD_aXD_a-\frac12\{G,X\}
 =\mathcal L_{q\ell}(X).
 \label{eq:v8-channel-generator-derivative}
\end{equation}
Но само семейство $\Phi_p$ не является точной полугруппой:
$\Phi_p\Phi_q\neq\Phi_{p+q}$. Для непрерывной марковской динамики
необходимы последовательность свежих состояний вспомогательных систем либо квантовый
шумовой носитель типа Фок-пространства. Текущая конструкция не выводит
такой поток и не выбирает функцию $p(t)$.

\begin{theorem}[Конечный ковариантный носитель одного шага]
Межсекторный канал допускает минимальную тринадцатимерную Стайнспринг-
дилатацию, чья двенадцатимерная скачковая часть совпадает с уже существующим
межсекторных стрелок представлением. Новая физическая частица для одного шага не
требуется. Однако вероятность шага и автономная непрерывная шумовая
динамика из родителя не следуют.
\end{theorem}

\begin{nogocriterion}[Запрет объявления одношагового канала временем]
Нельзя отождествлять параметр $p$ с физическим временем или скоростью без
правила композиции. Точная вполне положительная карта одного шага не
является сама по себе квантово-марковским потоком.
\end{nogocriterion}

\begin{protocol}
\begin{enumerate}
 \item размерность системного носителя: \textbf{21};
 \item независимых операторов скачка: \textbf{12};
 \item Чой/Краус-ранг при $0<p<1/6$: \textbf{13};
 \item допустимый интервал шага: \textbf{$[0,1/6]$};
 \item сохранение единицы и следа: \textbf{пройдено};
 \item сохранение концевой алгебры: \textbf{пройдено};
 \item калибровочная ковариантность среды: \textbf{пройдена};
 \item новый физический мультиплет: \textbf{не требуется};
 \item вероятность $p$: \textbf{не выведена};
 \item точная полугруппа $\Phi_p$: \textbf{нет};
 \item непрерывный шумовой/часовой носитель: \textbf{открыт}.
\end{enumerate}
\end{protocol}

Машинный сертификат сохранён в
\path{s2t_v8_minimal_covariant_stinespring_carrier_gate_results.json}.

\begin{thebibliography}{9}
\bibitem{v8-stinespring-haapasalo}
E.~Haapasalo,
\emph{The Choi--Jamio\l{}kowski isomorphism and covariant quantum channels},
Quantum Studies: Mathematics and Foundations \textbf{8} (2021), 1--30;
arXiv:1906.11442.
\bibitem{v8-stinespring-skeide}
M.~Skeide,
\emph{Spatial Markov semigroups admit Hudson--Parthasarathy dilations},
SIGMA \textbf{18} (2022), 071; arXiv:0809.3538.
\bibitem{v8-stinespring-parthasarathy}
K.~R.~Parthasarathy,
\emph{Quantum stochastic calculus and quantum Gaussian processes},
Indian Journal of Pure and Applied Mathematics \textbf{46} (2015),
781--807; arXiv:1408.5686.
\end{thebibliography}